\subsection{The risk of dying from COVID-19}
\totalscore{12}

%IEEE TRANSACTIONS ON ARTIFICIAL INTELLIGENCE, VOL. 02, NO. 1, FEB. 2021 1
%Simpson’s paradox in COVID-19 case fatality rates: a mediation analysis of age-related causal effects
%Julius von K\"ugelgen, Luigi Gresele, Bernhard Sch\"olkopf

In this exercise we investigate the influence of age and country on the risk of dying from a COVID-19 infection.
Figure~\ref{fig:covid_cfrs} displays data from the beginning of the COVID-19 pandemic in 2020, which started
with an outbreak in China followed a few weeks later by an outbreak in Italy.\footnote{Figure from
``Simpson's paradox in COVID-19 case fatality rates: a mediation analysis of age-related causal effects'',
Von K\"ugelgen, Gresele \& Sch\"olkopf, IEEE Transactions on Artificial Intelligence, 2(1), pp.\ 18--27, 2021.}
Figure~\ref{fig:covid_cfrs} (left) shows COVID-19 CFRs (case fatality rates,
that is, the proportion of confirmed COVID-19 cases that end with the death of the patient); 
for both countries, the CFRs are displayed for different age groups separately, and in addition the
total CFR for all age groups combined is shown.
Figure~\ref{fig:covid_cfrs} (right) shows for both countries how the confirmed COVID-19 cases
are distributed over the age groups.

Interestingly, while for all age groups, CFRs in Italy are lower than in China, the total CFR in Italy is
higher than that in China. This is an example of a phenomenon known as \emph{Simpson's paradox}.
%An explanation of this is that the distribution of ages differs between
%both countries. In particular, Italy recorded a much higher proportion of confirmed cases in older patients,
%as shown in Figure~\ref{fig:covid_cfrs} (right).

\begin{figure}
\includegraphics[width=\textwidth,page=2,trim={1.5cm 20cm 1cm 2cm},clip]{exercises/2005.07180v3.pdf}
  \caption{Left: COVID-19 CFRs for two countries (for different age groups, and total). Right: age distribution for confirmed COVID-19 cases in the two countries.\label{fig:covid_cfrs}}
\end{figure}

We decide to model the situation in the following way.
For each confirmed case (that is, a person that was tested positive in a certain time frame and reported by health officials of a certain country), we consider the following three variables:
\begin{center}\small\begin{tabular}{llll}
  Name & Symbol & Meaning \\
  \hline
  Country  & $X$ & Country in which the case is reported \\
  Age      & $Z$ & Age group to which the person belongs \\
  Fatality & $Y$ & Binary indicator of whether the person has deceased by the time of reporting
\end{tabular}\end{center}
We assume that the data-generating process can be modeled with a structural causal model with the following causal graph:
%\begin{figure}[!h]\centering
\begin{center}\begin{tikzpicture}
  \begin{scope}
    \node[ndout] (X) at (0,0) {$X$};
    \node[ndout] (Z) at (1.5,1) {$Z$};
    \node[ndout] (Y) at (3,0) {$Y$};
    \draw[huh] (X) to (Z);
    \draw[tuh] (Z) to (Y);
    \draw[tuh] (X) to (Y);
  \end{scope}
\end{tikzpicture}\end{center}
%  \caption{Country ($X$) and age ($Z$) have an influence on COVID-19 fatality ($Y$).}
%\end{figure}
\begin{enumerate}[label=\alph*)]
  \item Under what circumstances is it realistic to assume that the edges $Y \tuh X$ and $Y \tuh Z$ are both absent from the causal graph?
    \score{1}
    \answer{If the values of the variables ``country'' and ``age'' are measured \emph{before} ``fatality'' is measured.}
  \item Explain why it is sensible to assume that $X \tuh Z$ and $Z \tuh X$ are both absent from the causal graph.
    \score{1}
    \answer{Intervening on the country where someone lives does not change the age of that person, hence $X \tuh Z$ is absent.
    Similarly, intervening on the age of a person does not cause that person to move to another country, hence $Z \tuh X$ is absent.}
  \item Prove that if the edge $X \huh Z$ were absent from the causal graph, then country and age would be independent.
    \score{1}
    \answer{This follows from the global Markov property.
      If the bidirected edge $X \huh Z$ were absent from the graph $G(M)$ of the SCM $M$,
      (that is, if $G(M)$ were $X \tuh Y \hut Z$), 
      then $X \Perp^\sigma_G Z$, and hence (by the Markov property), $X \Indep_{P_M} Z$ where $P_M$ is the observational distribution associated to the SCM.}
  \item Give the technical term for what the edge $X \huh Z$ represents. 
    Explain with the help of Figure~\ref{fig:covid_cfrs} whether you think that the bidirected edge $X \huh Z$ should indeed be part of the model, or whether it could be safely omitted.
    \score{2}
    \answer{The bidirected edge $X \huh Z$ represents the confounding (spurious dependence, confounding bias) between ``country'' and ``age''.
      Figure~\ref{fig:covid_cfrs} (right) shows that for the confirmed cases, the age distribution in Italy is significantly different from that in China.
      In particular, in Italy there were many more confirmed cases among elderly persons. The observed dependence between age and country is explained by the
      bidirected edge $X \huh Z$, and would not be explained without it. Hence, the causal model would be incompatible with the data without the bidirected edge.
    }
%      Finally, bidirected edges $X \huh Y$ and $Z \huh Y$ are absent, which corresponds to assuming that there is no confounding between $X$ and $Y$, nor between $Z$ and $Y$.
%      Hence, the dependences between $(X,Z)$ and $Y$ are interpreted causally, that is, 
\end{enumerate}

%We now want to use the data and the causal model to answer the question:
%\emph{what would be the effect on fatality of changing country from China to Italy?}
One interesting quantity to consider is the
average total causal effect of country on fatality, defined as:
$$\rho := \Exp(Y \given \doit(X=1)) - \Exp(Y \given \doit(X=0))$$
%Note that here we interpret ``changing country'' as ``forcing all patients to move to another country''.

\begin{enumerate}[resume*]
  \item Give an identity that expresses this quantity purely in terms of the observed distribution $P(X,Y,Z)$.
    To simplify matters, you may assume that $P(X,Y,Z)$ is strictly positive.\\
    \emph{Hint: you may make use of the backdoor criterion.}
    \score{1}
\answer{
  The backdoor criterion states that,
  since $Y \Perp I_X \given X, Z$ and $Z \Perp I_X$,
    $$P(Y \given \doit(X)) = P(Y \given Z,X) \circ P(Z)$$
  under the absolute continuity assumption 
    $$P(Z) \otimes P(X) \ll P(Z,X).$$
  Writing this out, we get:
    $$\rho = \sum_z (\Exp(Y \given Z=z,X=1) - \Exp(Y \given Z=z,X=0)) P(Z=z) \text{ a.s.}$$
}
\end{enumerate}

Another question of interest is: for the Chinese case demographic, would the Italian health-care system have been better?
More specifically: what is the change in the risk of dying from COVID-19 when a randomly selected positive case from
China is forced to move to Italy.
We can model this by adding a copy $X'$ of $X$ (interpreted as the value of ``country'' measured later in time) and an intervention variable $I_{X'}$.
The causal graph then becomes:
\begin{center}\begin{tikzpicture}
    \node[ndout] (X) at (0,0) {$X$};
    \node[ndlat] (XX) at (1.5,0) {$X'$};
    \node[ndout] (Z) at (1.5,1) {$Z$};
    \node[ndout] (Y) at (3,0) {$Y$};
    \node[ndint] (IXX) at (1.5,-1.5) {$I_{X'}$};
    \draw[huh] (X) to (Z);
    \draw[tuh] (Z) to (Y);
    \draw[tuh] (X) to (XX);
    \draw[tuh] (XX) to (Y);
    \draw[tuh] (IXX) to (XX);
\end{tikzpicture}\end{center}
We assume that if we don't intervene on $X'$, its value is just $X' = X$.
In this setting, we can consider the quantity:\footnote{This quantity is also sometimes called the ``effect of treatment on the treated''.}
\begin{equation*}
  \tau := \Exp(Y \given X=0, \doit(X'=1)) - \Exp(Y \given X=0).
\end{equation*}
%Or is this the \emph{effect-of-treatment-on-the-treated}?
%Prove that the following \emph{mediation formula} holds:
\begin{enumerate}[resume*]
  \item Prove that $\tau$ can be expressed in terms of the observed distribution $P(X,Y,Z)$ as:
  \begin{equation*}\begin{split}
      \tau
      &= \sum_z (\Exp(Y \given Z=z, X=1) - \Exp(Y \given Z=z, X=0)) P(Z=z \given X=0) \text{ a.s.}
    \end{split}\end{equation*}
    To simplify matters, you may assume that $P(X,Y,Z)$ is strictly positive.\\
    \emph{Hint: the Markov kernel $P(X,X',Y,Z)$ is \emph{not} strictly positive.}
    \score{5}
\answer{
From the definition:
  $$\tau = \Exp(Y \given X=0, \doit(X'=1)) - \Exp(Y \given X=0).$$
The first term can be rewritten, using the do-calculus, as:
\begin{equation*}\begin{split}
  &\Exp(Y \given X=0, \doit(X'=1)) \\
  \text{[chain rule]} & = \sum_z \Exp(Y \given Z=z, X=0, \doit(X'=1)) P(Z=z \given X=0, \doit(X'=1)) \\
  \text{[$Z \Indep I_{X'} \given X$]} & = \sum_z \Exp(Y \given Z=z, X=0, \doit(X'=1)) P(Z=z \given X=0) \text{ a.s.} \\
  \text{[$Y \Indep X \given Z, X'$ under $\doit(X')$]} &= \sum_z \Exp(Y \given Z=z, \doit(X'=1)) P(Z=z \given X=0) \text{ a.s.} \\
  \text{[$Y \Indep I_{X'} \given Z, X'$]} &= \sum_z \Exp(Y \given Z=z, X'=1) P(Z=z \given X=0) \text{ a.s.} \\
  \text{[$X' = X$ a.s.]} &= \sum_z \Exp(Y \given Z=z, X=1) P(Z=z \given X=0) \text{ a.s.}
\end{split}\end{equation*}

For the second term we just use the chain rule:
$$\Exp(Y \given X=0) = \sum_z \Exp(Y \given Z=z, X=0) P(Z=z \given X=0).$$

Together, this proves the formula.
}
\points{The following alternative derivation that may look correct at first sight fails on positivity assumptions (because of the context-specific independence):
  \begin{equation*}\begin{split}
    &\Exp(Y \given X=0, \doit(X'=1)) \\
    &= \sum_z \Exp(Y \given Z=z, X=0, \doit(X'=1)) P(Z=z \given X=0, \doit(X'=1)) \\
    &= \sum_z \Exp(Y \given Z=z, X=0, X'=1) P(Z=z \given X=0) \text{ a.s.} \\
    &= \sum_z \Exp(Y \given Z=z, X'=1) P(Z=z \given X=0) \text{ a.s.} \\
    &= \sum_z \Exp(Y \given Z=z, X=1) P(Z=z \given X=0) \text{ a.s.}
  \end{split}\end{equation*}
  where we used the chain rule, $Y \Indep I_{X'} \given {X'}, Z, X$ and $Z \Indep I_{X'} \given X$, 
  $Y \Indep X \given {X'}, Z$ and that $X' = X$ a.s.

  I ended up subtracting 2 points if the positivity assumption was violated in the derivation.
}

\item How could you explain Simpson's paradox in this case?
    \score{1}
    \answer{While the risk of dying from COVID-19 is lower in Italy than in China for each age group (possibly reflecting the quality of the health-care system), it increases strongly with age.
    This dependence is so strong that in Italy, on average the CFR was higher than in China, as the relatively large fraction of elderly cases in Italy dominated the CFR.}
\end{enumerate}
